% Independent proof fragment; not automatically included in the manuscript.
% Basis and coordinate order matter for the reported determinant.
\begin{proposition}[One-product recovery and restricted minimality]
Retain the four local means
\[
 f=(-au,-bv,-cu,-dv)^T,\qquad
 (a,b,c,d)=(\sin\theta_0,\sin\theta_1,
           \sin\theta_2,\sin\theta_3),
\]
where $u=\cos x_0$, $v=\cos x_1$ and the $(u,v)$ image of the
input domain contains a nonempty open set. The four angles are
independently variable, and the kernel bandwidth $\gamma>0$ is fixed.
Among designs that retain these four means, append scalar features,
and use one input- and parameter-independent positive-definite
quadratic metric, the minimum number of appended scalar features
needed for generic kernel differential rank four is one. It is
attained by $\widetilde f=(f^T,f_0f_2)^T$ and the unit-trace metric
\[
 M=\frac{1}{7}M_0,\qquad
 M_0=\begin{pmatrix}
 2&1&0&0&0\\1&2&0&0&0\\0&0&1&0&0\\
 0&0&0&1&0\\0&0&0&0&1
 \end{pmatrix}.
\]
In contrast, appending only one pairwise product $f_if_j$ cannot
achieve rank four with a fixed positive diagonal metric.
\end{proposition}
\begin{proof}
With no appended feature, every fixed metric gives a symmetric
$2\times2$ coefficient Gram on the data basis $(u,v)$, so the kernel
differential rank is at most three. The same bound persists for any
fixed affine expansion of these four means, since its effective
quadratic metric leaves the data basis unchanged.

For the stated construction, $f_0f_2=ac\,u^2$. The functions
$1,u,v,u^2$ are linearly independent: a polynomial relation on an
open set in $(u,v)$ forces every coefficient to vanish. In
particular, on the actual domain $x\in(-1,1)^2$ the image contains
$(\cos1,1)^2$. Thus the coefficient-Gram theorem applies to
$h=(u,v,u^2)^T$, even though $u^2$ has a constant Fourier component.
The eigenvalues of $M_0$ are $3,1,1,1,1$, so $M\succ0$ and
$\operatorname{tr}M=1$. The unnormalized coefficient Gram is
\[
 G_0=\begin{pmatrix}
 2a^2+c^2&ab&0\\ab&2b^2+d^2&0\\0&0&a^2c^2
 \end{pmatrix},\qquad G=G_0/7.
\]
For the four entries in the order
$(G_{0,uu},G_{0,uv},G_{0,vv},G_{0,u^2u^2})$, the Jacobian in
$(a,b,c,d)$ is
\[
 \begin{pmatrix}
 4a&0&2c&0\\b&a&0&0\\0&4b&0&2d\\
 2ac^2&0&2a^2c&0
 \end{pmatrix},
\]
whose determinant is
\[
 -8a^2cd(2a^2-c^2).
\]
The angle Jacobian multiplies this by
$\prod_{q=0}^3\cos\theta_q$. At
$\theta_0=\theta_1=\theta_2=\theta_3=\pi/6$ its determinant is
$-9/128$. For the unit-trace metric the corresponding determinant is
\[
 -\frac{9}{128\,7^4}=-\frac{9}{307328}\ne0.
\]
Analyticity therefore gives generic rank four. No more than four
independent retained angles occur, giving the matching upper bound.
Zero appended coordinates cannot suffice and this one coordinate
does, establishing the stated restricted minimum.

For the diagonal-metric restriction, any one pairwise product has
the separable form $p(\theta)g(x)$ with
$g\in\{u^2,uv,v^2\}$. The resulting squared distances depend on
at most the three scalar coefficients
\[
 w_a a^2+w_c c^2,\qquad w_b b^2+w_d d^2,\qquad
 w_p p(\theta)^2.
\]
Their differential rank is at most three, for every fixed choice of
positive diagonal weights.
\end{proof}

\paragraph{Scope of the minimum.}
The minimum counts appended scalar outputs while retaining the
original four means and using a fixed quadratic metric. It is not a
minimum over arbitrary replacement representations, input-dependent
or parameter-dependent metrics, trainable postprocessing, or other
kernel families. The diagonal-metric obstruction is specifically
for one \emph{pairwise-product coordinate}, not an arbitrary scalar
coordinate. Indeed,
\[
 g=f_0f_2+f_1f_3=ac\,u^2+bd\,v^2
\]
is one scalar whose coefficient Gram under positive diagonal weights
already exposes
\[
 w_a a^2+w_c c^2,\quad w_g a^2c^2,\quad
 w_b b^2+w_d d^2,\quad w_g b^2d^2.
\]
The two $(a,c)$ and $(b,d)$ Jacobian blocks are generically
nonsingular. Hence this different one-scalar lift attains rank four
even with fixed positive diagonal weighting.

\paragraph{Generic rank does not imply global uniqueness.}
The displayed determinant vanishes at exceptional points; for
example all distance first derivatives vanish when all four sines
are zero. Even away from the exceptional set, global uniqueness
fails. The distinct positive amplitude tuples
\[
 (a,b,c,d)=(1/2,1/4,1/2,1/2),\qquad
 (a,b,c,d)=(1/(2\sqrt2),\sqrt2/4,1/\sqrt2,\sqrt2/4)
\]
are both regular and give the same ordered four Gram entries
$(3/4,1/8,3/8,1/16)$ for $G_0$. Sine periodicity and sign ambiguities
provide further nonuniqueness. The result identifies the four
retained angles locally at generic points, not the diagonal-tail
symbols, and makes no conditioning, shot-efficiency, or predictive
performance claim.

\paragraph{No additional ideal measurement information.}
In the stated circuit family the state before the diagonal tail is
a product, and the tail commutes with every Z string. Consequently
$\langle Z_0Z_2\rangle=\langle Z_0\rangle\langle Z_2\rangle=f_0f_2$.
The appended ideal feature is therefore a nonlinear classical
function of the four local means. It can also be estimated as a
parity using the same Z-basis outcomes. Direct parity averages and
products of local sample means are different finite-shot estimators;
the ideal identity alone does not equate their precision. General
correlated states need not satisfy this identity.

\paragraph{Basis convention for exact certificates.}
The determinants above use $h=(u,v,u^2)$. If a strictly nonconstant
Fourier-mode basis $(\cos x_0,\cos x_1,\cos2x_0)$ is used instead,
the last coefficient is $ac/2$ and the last feature has offset $ac/2$.
The corresponding Gram-entry Jacobian minor is smaller by a factor
of four: its angular value is $-9/512$ before trace normalization,
or $-9/1229312$ after normalization. This basis change does not
alter the kernel or its rank.
